\subsection{Réduction causale de la mise à jour bilatérale}
\label{lib04:reduccion-causal}

La mise à jour unitaire précédente réalise le transport de l’état et de sa
mémoire dans un même espace. L’observation de la première composante admet
une description équivalente qui conserve explicitement la contribution du
registre. Cette description s’obtient à partir du transport déjà construit, avec les
mêmes coefficients et la même orientation ; la mémoire éliminée de la liste
des variables réapparaît comme dépendance à l’égard des états précédents.

Fixons les opérateurs \(A,D\) de \eqref{lib04:ad} et une réalisation \(H\oplus H\). Pour \(n\geq0\), soient
\begin{equation}
 v_{n+1}=Av_n-D^*m_n,\qquad
 m_{n+1}=Dv_n+A^*m_n.
 \label{lib04:recurrencia-memoria}
\end{equation}
La mise à jour conjointe est \(\mathscr R_a\) et conserve \(\|v_n\|^2+\|m_n\|^2\).

\begin{proposition}[Élimination exacte du registre auxiliaire]
\label{lib04:prop-memoria-reducida}
Pour tout \(n\geq0\), la somme vide étant interprétée comme zéro,
\begin{align}
 m_n&=(A^*)^n m_0+
       \sum_{k=0}^{n-1}(A^*)^{n-1-k}Dv_k,
       \label{lib04:memoria-explicita}\\
 v_{n+1}&=Av_n-D^*(A^*)^n m_0
       +\sum_{k=0}^{n-1}\mathscr K_{n-1-k}v_k,
       \label{lib04:evolucion-reducida}\\
 \mathscr K_j&=-D^*(A^*)^jD\quad(j\geq0).
       \label{lib04:nucleo-retardo}
\end{align}
L’équation réduite, avec \eqref{lib04:memoria-explicita}, est
équivalente à la récurrence conjointe pour les mêmes données \((v_0,m_0)\).
\end{proposition}
\begin{proof}
La première identité vaut pour \(n=0\). Si elle vaut pour \(n\), la seconde équation de \eqref{lib04:recurrencia-memoria} fournit
\[
 m_{n+1}=(A^*)^{n+1}m_0+
 \sum_{k=0}^{n-1}(A^*)^{n-k}Dv_k+Dv_n,
\]
qui est l’identité dans \(n+1\). La substituer, à la profondeur \(n\), dans
la première équation démontre \eqref{lib04:evolucion-reducida}.
Réciproquement, si \(v_n\) satisfait cette dernière et si l’on définit \(m_n\)
par \eqref{lib04:memoria-explicita}, la même différence de sommes prouve
\(m_{n+1}=Dv_n+A^*m_n\) ; l’autre équation se récupère par substitution.
\end{proof}

Le noyau \(\mathscr K_j\) décrit un transfert vers le registre, \(j\) mises à jour à l'intérieur de celui-ci et une rentrée dans la composante lue. Le terme \(-D^*(A^*)^nm_0\) conserve la mémoire initiale. Sa suppression correspond au cas particulier \(m_0=0\), et non à une identité du transport général. Dans \eqref{lib04:evolucion-reducida} interviennent uniquement l'état actuel et les états antérieurs. La causalité de cette réduction provient de l'ordre de la récurrence conjointe.

La normalisation du bilan fournit en outre une estimation explicite.
Comme \(Q_-=aI-U\), on a
\(\|A\|^2\leq r_a=(a+1)^3/((a+1)^3+a^3)<1\), tandis que
\(D^*D\leq I\). Par conséquent,
\begin{equation}
 \|\mathscr K_j\|\leq r_a^{j/2},\qquad
 \sum_{j\geq0}\|\mathscr K_j\|
 \leq\frac{1}{1-\sqrt{r_a}}.
 \label{lib04:cota-nucleo-retardo}
\end{equation}
Cette borne contrôle le poids de la mémoire dans l’équation réduite ; la norme
de l’état conjoint reste conservée par la mise à jour unitaire.

La proposition est la réalisation temporelle de la réduction de Schur de
\cref{mem:prop-schur}. Dans les séries formelles
\(V(z)=\sum_{n\geq0}v_nz^n\) et
\(M(z)=\sum_{n\geq0}m_nz^n\), les récurrences donnent
\(V=v_0+zAV-zD^*M\) et \(M=m_0+zDV+zA^*M\). En éliminant \(M\),
\begin{equation}
 V(z)=\bigl[I-zA+z^2D^*(I-zA^*)^{-1}D\bigr]^{-1}
       \bigl[v_0-zD^*(I-zA^*)^{-1}m_0\bigr].
 \label{lib04:schur-bilateral}
\end{equation}
Les inverses formels existent car leurs termes constants sont des identités ;
les séries des états et la résolvante conjointe convergent également pour
\(|z|<1\). La source initiale et l’opérateur effectif se transforment ensemble.
Un lecteur qui agit également sur \(m_n\) conserve sa contribution au moyen de
\eqref{lib04:memoria-explicita}, conformément à \eqref{mem:schur-lector}.

La réentrée décrite ici actualise un registre entrant. La politique
d’archivage qui suit conserve en revanche chaque complément dans une coordonnée
distincte. Les deux opérations possèdent leur propre loi de conservation. Le noyau
\(\mathscr K_j\), les propagateurs de Green et une force dimensionnelle sont
des objets de types différents : le premier élimine une composante d’état,
les seconds résolvent des équations de source avec des conditions de support, et
la force ajoute une réalisation dynamique et ses unités.
